% TEMPLATE-MILC-v3.tex - plantilla maestra UNICA de las guias MILC v3.
%
% La usan las cuatro familias de guias, cada una con su driver:
%   · Grados 6-11          -> scripts/build-guias-g11.py
%   · Bebras               -> scripts/build-guias-bebras.py
%   · Semillero            -> scripts/build-guias-semillero.py
%   · Territorio interior  -> scripts/build-guias-territorio-interior.py
%
% Antes habia tres copias casi identicas (~400 lineas iguales, 6 distintas):
% cada mejora habia que aplicarla tres veces, y en la practica no se hacia
% (el motor de recursos vivia solo en dos de ellas). Lo que varia entre
% familias esta parametrizado en 6 placeholders que cada driver llena
% (se nombran aqui SIN su sintaxis real: el builder reemplaza por texto plano,
% tambien dentro de los comentarios, y un valor multilinea rompe el preambulo):
%
%   HEADER_IZQ        encabezado izquierdo de pagina
%   HEADER_DER        encabezado derecho de pagina
%   PORTADA_ETIQUETA  rotulo sobre el numero grande (GRADO/MOMENTO/MODULO)
%   PORTADA_SERIE     linea de serie de la portada
%   PORTADA_ID        identificador de la guia en la portada
%   PORTADA_CIRCULO   el circulito de grado (°), o vacio si no aplica
%   PDF_TITULO        titulo en texto plano (sin \textbf ni comillas TeX) para
%                     los metadatos del PDF (/Title); lo produce lib_xelatex.texto_pdf
%
% Composicion (auditoria 2026-09-03): las bandas de estacion piden espacio
% (\Needspace*) y prohiben el salto justo despues (\nopagebreak), asi no quedan
% huerfanas al pie; las cajas largas (softbox, drawbox, infoband, puentebox) son
% `breakable`, asi una caja de media pagina ya no arrastra todo a la siguiente
% dejando la anterior al 40 %. Todo texto sobre mostaza o verde va en milcNegro
% (blanco sobre #E5B400 da 1,93:1 y sobre #58A923 2,95:1; ninguno pasa WCAG AA).
% El resto de claves las documenta content/guias/_SCHEMA.md. El script valida
% que no queden placeholders sin reemplazar antes de escribir el archivo final.

% Etiquetado PDF (latex-lab): /Lang, estructura y texto alternativo de las
% figuras, para lectores de pantalla. Debe ir ANTES de \documentclass.
\DocumentMetadata{lang=es-CO,tagging=on}
\documentclass[11pt,letterpaper]{article}
% headheight=24pt: la cabecera derecha lleva dos lineas (identidad de la guia
% y estacion actual). headsep se acorta en la misma medida para que el bloque
% de texto quede exactamente donde estaba con headheight=12pt.
\usepackage[margin=2.0cm,headheight=24pt,headsep=13pt]{geometry}
\usepackage[table]{xcolor}
\usepackage{fontspec}
\usepackage[spanish]{babel}
\usepackage{tikz}
% Librerías TikZ para los diagramas de las guías (bloque `recursos:`):
%  · babel  → babel en español vuelve activos caracteres como «>», y TikZ los usa
%             en sus opciones (>=stealth, ->). Sin esta librería, cualquier
%             diagrama con flechas revienta con «\language@active@arg> has an
%             extra }». Debe ir ANTES de que se dibuje nada.
%  · positioning        → sintaxis «below left=2pt and 3pt».
%  · decorations.pathreplacing → llaves ({brace}) para señalar rangos.
\usetikzlibrary{babel,positioning,decorations.pathreplacing}
\usepackage{graphicx}
% Circuitos: el trabajo de electrónica se hace hoy con micro:bit y sus sensores
% integrados (MakeCode, sin cableado), así que no se necesitan esquemáticos.
% Si algún día entran montajes con protoboard, basta descomentar esta línea:
% los diagramas .tex/.tikz declarados en `recursos:` ya se incrustan solos.
% \usepackage{circuitikz}
\usepackage[most]{tcolorbox}
\usepackage{tabularx}
\usepackage{array}
\usepackage{fancyhdr}
\usepackage{titlesec}
\usepackage[document]{ragged2e}  % bandera a la derecha en todo el cuerpo
\usepackage{enumitem}
\usepackage{microtype}
\usepackage{etoolbox}
\usepackage{needspace}
% hyperref al final del preambulo: metadatos del PDF (titulo, autor, idioma,
% familia), marcadores por estacion (\pdfbookmark) y enlaces sin recuadro.
\usepackage[hidelinks,
  pdftitle={La balanza y el peso — medir como acto político},
  pdfauthor={PhD. Álvaro Cárdenas Orozco},
  pdfsubject={Tecnología e Informática · Grado 9 · Guía 3 · MILC v3},
  pdflang={es-CO},
  bookmarksopen=true,bookmarksnumbered=false]{hyperref}

% Helvetica Neue no trae la flecha U+2192, asi que cada «→» escrito literalmente
% en el YAML se compilaba a NADA, en silencio (462 flechas en 61 guias: rutas del
% tipo «paleta → tipografia → lockup» salian rotas). Mapeamos el caracter a su
% version matematica, que si tiene glifo. Asi el autor puede seguir escribiendo
% «→» comodamente en el YAML.
\usepackage{newunicodechar}
\newunicodechar{→}{\ensuremath{\rightarrow}}
% Mismo problema con los simbolos de marca y vinetas: el «✓» de «una palomita ✓»
% salia como cuadrito vacio en el PDF de todas las guias que lo usaban.
\usepackage{pifont}
\newunicodechar{✓}{\ding{51}}
\newunicodechar{✗}{\ding{55}}
\newunicodechar{✔}{\ding{52}}
\newunicodechar{•}{\textbullet}
\newunicodechar{≥}{\ensuremath{\geq}}
\newunicodechar{≤}{\ensuremath{\leq}}

\setmainfont{Helvetica Neue}
\setsansfont{Helvetica Neue}
\setlength{\parindent}{0pt}
\setlength{\parskip}{6pt}
% Sin particion de palabras: en 6.º-7.º una palabra cortada por guion al final
% de linea («Comuni-cacion») frena la lectura mucho mas que una linea corta.
\hyphenpenalty=10000
\exhyphenpenalty=10000
\pretolerance=2000
\tolerance=3000
\setlength{\emergencystretch}{3em}
\linespread{1.10}
\renewcommand{\arraystretch}{1.25}
\setlist{leftmargin=5mm,itemsep=1mm,topsep=1mm}
\newcolumntype{Y}{>{\RaggedRight\arraybackslash}X}

\definecolor{milcMagenta}{HTML}{E6007E}
\definecolor{milcVino}{HTML}{5A0038}
\definecolor{milcTurquesa}{HTML}{0093A5}
\definecolor{milcVerde}{HTML}{58A923}
\definecolor{milcMostaza}{HTML}{E5B400}
\definecolor{milcOcre}{HTML}{B86600}
\definecolor{milcNegro}{HTML}{241816}
\definecolor{milcGris}{HTML}{F7F7F4}
\definecolor{milcLinea}{HTML}{E8E2DD}
\definecolor{milcGradePrimary}{HTML}{7C3AED}
\definecolor{milcGradeDark}{HTML}{4C1D95}
\definecolor{milcGradeSoft}{HTML}{EDE9FE}
\definecolor{milcGradeAccent}{HTML}{CFE7FF}
\definecolor{milcGradeText}{HTML}{FFFFFF}
\color{milcNegro}

\pagestyle{fancy}
\fancyhf{}
% Cabecera de dos lineas: identidad de la guia arriba y, a la derecha y en
% gris, la estacion en curso (\markright desde \milcestacion). Antes de la
% primera estacion la segunda linea va vacia; el \strut mantiene la altura
% para que la linea de arriba no baile entre paginas.
\lhead{\small Tecnología e Informática · Grado 9\\[-1pt]{\footnotesize\strut}}
\rhead{\small Guía 3 · MILC v3\\[-1pt]{\footnotesize\strut\color{milcNegro!65}\rightmark}}
\cfoot{\small\thepage}
\renewcommand{\headrulewidth}{0pt}

% Sobre mostaza (#E5B400) y verde (#58A923) el texto va en milcNegro; sobre el
% resto de la paleta, en blanco. Se usa dentro de fonttitle/fontupper, donde un
% \color posterior manda sobre coltitle.
\newcommand{\milctintasobre}[1]{%
  \ifboolexpr{test{\ifstrequal{#1}{milcMostaza}} or test{\ifstrequal{#1}{milcVerde}}}%
    {\color{milcNegro}}{\color{white}}}

\newtcolorbox{titlebox}[1]{
  enhanced,colback=#1,colframe=#1,arc=7mm,boxrule=0pt,
  left=7mm,right=7mm,top=3mm,bottom=3mm,
  before skip=8pt,after skip=8pt,
  fontupper=\sffamily\bfseries\Large\color{white}
}

\newtcolorbox{softbox}[3]{
  enhanced,breakable,pad at break=3mm,
  colback=#2,colframe=#1,borderline west={4pt}{0pt}{#1},
  arc=2mm,boxrule=.4pt,left=4mm,right=4mm,top=3mm,bottom=3mm,
  before skip=6pt,after skip=8pt,title={#3},coltitle=white,
  colbacktitle=#1,fonttitle=\sffamily\bfseries\milctintasobre{#1},fontupper=\RaggedRight,
  attach boxed title to top left={xshift=3mm,yshift=-2mm},
  boxed title style={arc=2mm, boxrule=0pt}
}

\newtcolorbox{drawbox}[2]{
  enhanced,breakable,pad at break=4mm,
  colback=white,colframe=#1,arc=4mm,boxrule=1pt,
  left=5mm,right=5mm,top=4mm,bottom=4mm,before skip=8pt,after skip=8pt,
  title={#2},coltitle=white,colbacktitle=#1,fonttitle=\sffamily\bfseries\milctintasobre{#1},
  fontupper=\RaggedRight,
  attach boxed title to top left={xshift=4mm,yshift=-2mm},
  boxed title style={arc=3mm, boxrule=0pt}
}

\newtcolorbox{infoband}[1]{
  enhanced,breakable,pad at break=2mm,colback=#1!8,colframe=#1!50,arc=2mm,boxrule=.5pt,
  left=4mm,right=4mm,top=2mm,bottom=2mm,before skip=4pt,after skip=4pt,
  fontupper=\small\RaggedRight
}

% Puente narrativo entre secciones (contrato editorial Conéctate).
% Da continuidad: "Acabas de X. Ahora vas a Y porque Z."
% Visualmente: banda izquierda en color del grado, texto en itálica pequeña.
\newtcolorbox{puentebox}{
  enhanced,breakable,pad at break=2mm,colback=milcGradeSoft,colframe=milcGradeDark,
  borderline west={3pt}{0pt}{milcGradeDark},
  arc=0mm,boxrule=0pt,left=5mm,right=4mm,top=2.5mm,bottom=2.5mm,
  before skip=4pt,after skip=8pt,
  fontupper=\itshape\small\color{milcNegro}\RaggedRight
}
% La flecha va en modo matematico a proposito: Helvetica Neue NO tiene el glifo
% U+2192, asi que \textrightarrow se compilaba a NADA («Missing character: There
% is no → in font Helvetica Neue») y el puente salia sin flecha. En modo
% matematico la flecha viene de la fuente matematica y si se dibuja.
\newcommand{\puente}[1]{%
  \begin{puentebox}%
    \textcolor{milcGradeDark}{$\rightarrow$}\hspace{2mm}#1%
  \end{puentebox}%
}

\newcommand{\linead}[1]{\noindent\color{milcLinea}\rule{#1}{0.5pt}\color{milcNegro}}
\newcommand{\respuesta}[1]{\par\vspace{2mm}\foreach \n in {1,...,#1}{\noindent\color{milcLinea}\rule{\linewidth}{0.5pt}\par\vspace{5mm}}\color{milcNegro}}
\newcommand{\checkbox}{\textcolor{milcGradeDark}{\fbox{\phantom{x}}}\hspace{5pt}}

% ─── Listas aireadas para las guías socioemocionales ────────────────────
% Componer las verificaciones como «\checkbox texto\\» deja el texto que salta
% de línea debajo del cuadrito y apelmaza el bloque. Estos entornos alinean la
% continuación y airean los ítems. Los usa Territorio Interior; cualquier
% familia puede hacerlo.
\newlist{checklista}{itemize}{1}
\setlist[checklista]{
  label=\textcolor{milcGradeDark}{\fbox{\phantom{x}}},
  leftmargin=9mm, labelsep=3.2mm, itemsep=2.4mm,
  topsep=2.5mm, parsep=0pt, partopsep=0pt, before=\RaggedRight
}
\newlist{pasoslista}{enumerate}{1}
\setlist[pasoslista]{
  label=\textcolor{milcGradeDark}{\sffamily\bfseries\arabic*.},
  leftmargin=8mm, labelsep=3mm, itemsep=2.8mm,
  topsep=1mm, parsep=0pt, partopsep=0pt, before=\RaggedRight
}

\newcommand{\phasecard}[4]{
\begin{tcolorbox}[enhanced,colback=#3,colframe=#2,arc=3mm,boxrule=.5pt,height=3.05cm,valign=center,halign=center,left=1.5mm,right=1.5mm,top=2mm,bottom=2mm]
{\sffamily\bfseries\fontsize{22}{24}\selectfont\textcolor{#2}{#1}}\par
{\sffamily\bfseries\small #4}
\end{tcolorbox}
}

% ── Piezas del rediseno 2026 (legibilidad 6.º-11.º) ───────────────────
% Banda de estacion: reemplaza el titlebox macizo. Titulo a la izquierda,
% dato practico (tiempo/modalidad) a la derecha, en una sola linea.
% Nunca huerfana: pide 9 lineas libres (o salta de pagina rellenando con
% \vfill) y prohibe el salto justo despues de la banda. Nueve y no seis:
% la caja partible que sigue necesita ~80pt (titulo adosado, relleno y dos
% lineas) para arrancar en la misma pagina; con menos, tcolorbox la manda
% entera a la siguiente y la banda se queda sola. El penalty 10000 va
% en `after`, ANTES del espacio posterior: un `after skip` (aunque sea 0pt)
% es un \vskip tras la caja, y ese glue es un punto de corte legal que un
% \nopagebreak escrito despues de \end{tcolorbox} ya no cierra. Cada banda
% deja marcador PDF y actualiza la cabecera derecha.
\newcounter{milcestacion}
\newcommand{\milcestacion}[2]{%
\Needspace*{9\baselineskip}%
\stepcounter{milcestacion}%
\pdfbookmark[1]{#2}{milcestacion\themilcestacion}%
\markright{#2}%
\begin{tcolorbox}[enhanced,colback=#1,colframe=#1,arc=2mm,boxrule=0pt,
  left=5mm,right=5mm,top=2.6mm,bottom=2.6mm,before skip=11pt,
  after={\par\nopagebreak\vskip7pt}]
{\sffamily\bfseries\fontsize{15}{18}\selectfont\milctintasobre{#1}#2}
\end{tcolorbox}}

% Tarjeta de paso numerada: sustituye las tablas «Pilar / Aplicacion», que
% eran formato de consulta docente, no de estudiante. El numero va en un
% circulo sobre el borde para que la secuencia se lea de un vistazo.
\newcommand{\milcpaso}[3]{%
\Needspace{4\baselineskip}%
\begin{tcolorbox}[enhanced,colback=white,colframe=milcLinea,arc=1.5mm,boxrule=.9pt,
  left=14mm,right=4mm,top=3mm,bottom=3mm,before skip=4pt,after skip=4pt,
  fontupper=\RaggedRight,
  overlay={\node[anchor=north west,circle,fill=#2,text=white,inner sep=0pt,
    minimum size=7.4mm,font=\sffamily\bfseries\small]
    at ([xshift=3.6mm,yshift=-3.2mm]frame.north west) {#1};}]
#3
\end{tcolorbox}}

% Casilla marcable de tamano util con lapiz (la anterior era \fbox{\phantom{x}},
% de alto variable segun el texto que la seguia).
\newcommand{\milcmarca}{\framebox[3.8mm]{\rule{0pt}{3.4mm}}}

% Fila marcable de lista suelta (checks de la estacion 1).
\newcommand{\milccheckrow}[1]{%
\par\vspace{1.8mm}\noindent
\makebox[6.4mm][l]{\raisebox{-.55mm}{\textcolor{milcNegro}{\milcmarca}}}%
\parbox[t]{\dimexpr\linewidth-6.4mm\relax}{\RaggedRight #1}\par}

% Celda marcable dentro de la rubrica (tabla de autoevaluacion). Misma
% sangria francesa, pero pensada para vivir dentro de una columna X.
\newcommand{\milcopcion}[1]{%
\makebox[5.2mm][l]{\raisebox{-.55mm}{\textcolor{milcNegro}{\milcmarca}}}%
\parbox[t]{\dimexpr\linewidth-5.2mm\relax}{\RaggedRight #1}}

% ── Recursos visuales de la guía ──────────────────────────────────────
% Declarados en el bloque `recursos:` del YAML y emitidos por el builder.
% \guiaFigura   → imagen rasterizada o PDF (foto, carta celeste, montaje).
% \guiaDiagrama → fuente TikZ/circuitikz (.tex) incrustada como vector:
%                 es la vía para circuitos y esquemáticos editables.
% [#1] = texto alternativo (alt) · #2 = ruta relativa al .tex · #3 = pie
% El alt llega al PDF etiquetado (/Alt) y lo lee el lector de pantalla.
\newcommand{\guiaFigura}[3][]{%
\begin{center}
\includegraphics[width=0.88\linewidth,height=0.38\textheight,keepaspectratio,alt={#1}]{#2}\\[1.5mm]
{\footnotesize\itshape\textcolor{milcNegro!75}{#3}}
\end{center}
}

\newcommand{\guiaDiagrama}[2]{%
\begin{center}
\input{#1}\\[1.5mm]
{\footnotesize\itshape\textcolor{milcNegro!75}{#2}}
\end{center}
}

\begin{document}
\thispagestyle{empty}

% ── Portada ───────────────────────────────────────────────────────────
% Antes: un rectangulo de color a pagina completa. Se veia bien en pantalla,
% pero en la fotocopiadora del colegio se comia el toner y salia gris sucio.
% Ahora la banda ocupa el tercio superior y el resto va en blanco.
\begin{tikzpicture}[remember picture,overlay]
  \path[fill=milcGradePrimary]
    (current page.north west) rectangle ([yshift=-10.6cm]current page.north east);

  \node[anchor=north west,text=milcGradeText] at ([xshift=2.0cm,yshift=-1.55cm]current page.north west)
    {\sffamily\bfseries\fontsize{12}{14}\selectfont GRADO};
  \node[anchor=north west,text=milcGradeText,opacity=.85] at ([xshift=2.0cm,yshift=-2.35cm]current page.north west)
    {\sffamily\bfseries\fontsize{8.5}{10}\selectfont SERIE GUÍAS · TECNOLOGÍA E INFORMÁTICA};
  \node[anchor=north east,text=milcGradeText,opacity=.85,align=right] at ([xshift=-2.0cm,yshift=-1.55cm]current page.north east)
    {\sffamily\bfseries\fontsize{9}{12}\selectfont I.E. Sor María Juliana\\Cartago, Valle del Cauca};

  \node[anchor=west,text=milcGradeText] (gradonum) at ([xshift=1.85cm,yshift=-7.1cm]current page.north west)
    {\sffamily\bfseries\fontsize{112}{112}\selectfont 9};
  % El circulito de grado (°) solo aplica a las guias de grado («11°»); en
  % Bebras y Semillero el numero grande es un momento/modulo, no un grado.
  % Se posiciona relativo al nodo `gradonum`, no en coordenadas absolutas.
  \draw[milcGradeText,line width=1.6pt]
    ([xshift=2.0mm,yshift=-4.5mm]gradonum.north east) circle (.15cm);

  \node[anchor=west,text=milcGradeText,text width=11.4cm,align=left] at ([xshift=1.5cm,yshift=0.5cm]gradonum.east)
    {
     {\sffamily\bfseries\fontsize{9}{11}\selectfont\MakeUppercase{Período 1 · Historia de la técnica}}\\[3mm]
     {\sffamily\bfseries\fontsize{21}{25}\selectfont\setlength{\baselineskip}{27pt}La balanza\\[4pt]y el peso ---\\[4pt]medir es decidir}\\[3.5mm]
     {\sffamily\bfseries\fontsize{15}{18}\selectfont Guía 3-9-TIC}
    };
\end{tikzpicture}

\vspace*{9.4cm}
\vfill

{\sffamily\bfseries\fontsize{8}{10}\selectfont\textcolor{milcGradeDark}{LO QUE TE LLEVAS HOY}}

\begin{tcolorbox}[enhanced,colback=white,colframe=milcNegro,arc=2mm,boxrule=1.6pt,
  left=5mm,right=5mm,top=4mm,bottom=4mm,before skip=3pt,after skip=12pt,
  fontupper=\RaggedRight\fontsize{11}{15.5}\selectfont]
Una balanza casera que se equilibre en vacío y permita comparar dos objetos, con su dibujo o foto y las piezas rotuladas. La acompaña una reflexión de una página sobre un sistema de medición que rige tu vida: qué mide, con qué patrón, quién definió ese patrón, qué decide a partir del resultado y a quién le conviene que se mida así.
\end{tcolorbox}

\begin{tcolorbox}[enhanced,colback=milcGris,colframe=milcGris,arc=2mm,boxrule=0pt,
  left=5mm,right=5mm,top=3.5mm,bottom=3.5mm,after skip=12pt]
{\sffamily\bfseries\fontsize{8}{10}\selectfont\textcolor{milcNegro!55}{ESTA GUÍA ES DE}}\\[3mm]
\begin{tabularx}{\linewidth}{X p{3.2cm} p{3.2cm}}
\linead{\linewidth} & \linead{3.0cm} & \linead{3.0cm}\\[-1mm]
{\fontsize{8.5}{10}\selectfont\textcolor{milcNegro!55}{Nombre completo}} &
{\fontsize{8.5}{10}\selectfont\textcolor{milcNegro!55}{Grupo}} &
{\fontsize{8.5}{10}\selectfont\textcolor{milcNegro!55}{Fecha}}\\
\end{tabularx}
\end{tcolorbox}

\vfill

\noindent\textcolor{milcLinea}{\rule{\linewidth}{1pt}}\\[2mm]
{\sffamily\bfseries\fontsize{9.5}{12}\selectfont Autor: Dr. Álvaro Cárdenas Orozco}\\[1mm]
{\sffamily\fontsize{8.5}{11}\selectfont\textcolor{milcNegro!55}{Metodología MILC · apertura ancestral · pensamiento computacional · triángulo Dussel–estoicismo–Floridi}}

\newpage

\section*{Apertura: La balanza y el peso --- medir como acto político}

\begin{softbox}{milcGradeDark}{milcGradeSoft}{Saber ancestral}
Entre 1949 y 1953, el sociólogo Orlando Fals Borda vivió en la vereda de Saucío, en Chocontá, Cundinamarca. Conviene decirlo desde el principio: es Cundinamarca, no el Valle. Allí encontró que la gente se medía de dos maneras a la vez. Una era lo que él llamó la \textbf{escala agrícola}: una escalera invisible con los propietarios arriba y los concertados sin tierra abajo. Hasta los muebles y los zapatos servían de indicador de posición. Fals Borda anotó algo incómodo: el campesino ahorraba sobre todo para sostener su prestigio, y ese prestigio se medía por su capacidad de corresponder invitaciones. La otra manera de medir era el convite y el brazo prestado. Ahí uno valía por lo que ponía cuando a otro le hacía falta. La \textbf{cara de exclusión}: no hay que idealizar la vereda. Fals Borda registró que en Saucío no había cooperativas ni reuniones para discutir problemas comunes. Vio desaparecer la segunda medida y vio la vereda recuperarse cuando volvieron a hacer algo juntos: construir la escuela.\par\smallskip{\footnotesize\textcolor{milcNegro!70}{\textbf{Fuente.} Fals Borda, O. (2017). \emph{Campesinos de los Andes y otros escritos antológicos}. Universidad Nacional de Colombia. (Obra original publicada en 1955)}}
\end{softbox}

\begin{softbox}{milcTurquesa}{milcGris}{Saber contemporáneo}
Una balanza tiene tres partes y ninguna sobra. El \textbf{punto de apoyo}, donde el sistema se equilibra. Los \textbf{dos brazos iguales}, uno para lo desconocido y otro para el patrón. Y el \textbf{patrón conocido}, que son las pesas calibradas. Ese diseño tan simple sostiene la confianza de un mercado entero. Si los dos brazos se equilibran con un patrón que ambas partes aceptan, las dos pueden confiar en el resultado. Sin patrón aceptado no hay medición: hay una opinión con aparato. Por eso medir nunca fue un asunto técnico solamente. Alguien decide qué se mide, con qué unidad y qué pasa contigo según el número que salga. Tu promedio, tus horas de pantalla, tus seguidores y las calorías de un producto son patrones que alguien definió. La pregunta útil no es si el número es exacto. Es quién eligió medir eso y no otra cosa.
\end{softbox}

\begin{softbox}{milcMostaza}{milcGris}{Pregunta-puente}
\textit{En Saucío convivían dos medidas del valor y cada una repartía prestigio de un modo distinto. ¿Qué dos medidas conviven en tu curso, y cuál pesa más a la hora de decidir quién cuenta?}
\end{softbox}

\begin{softbox}{milcVerde}{milcGris}{Saber y saber hacer}
Al terminar podrás: (1) \textbf{identificar} cinco sistemas de medición que regulan tu vida y quién definió cada patrón; (2) \textbf{explicar} las tres partes de una balanza y por qué el patrón compartido es lo que sostiene la confianza; (3) \textbf{crear} una balanza casera que funcione y una reflexión sobre a quién le conviene un sistema de medición concreto.
\end{softbox}



\small
\begin{tabularx}{\textwidth}{>{\bfseries}p{3.0cm}Y}
\rowcolor{milcGradePrimary!12}Autor & Dr. Álvaro Cárdenas Orozco\\
DBA & Análisis crítico de la historia de la tecnología y su impacto social (MEN, T\&I)\\
Referentes & Dussel (técnica situada) · Estoicismo (téchne del cuerpo) · Floridi (infoesfera)\\
Producto final & Una balanza casera que se equilibre en vacío y permita comparar dos objetos, con su dibujo o foto y las piezas rotuladas. La acompaña una reflexión de una página sobre un sistema de medición que rige tu vida: qué mide, con qué patrón, quién definió ese patrón, qué decide a partir del resultado y a quién le conviene que se mida así.\\
\end{tabularx}
\normalsize

\puente{En la sesión 1 definiste qué es técnica y en la 2 estudiaste máquinas simples del campo. La balanza es la máquina simple más antigua que sigue decidiendo cosas todos los días. De aquí sale la mirada crítica que vas a usar el resto del año.}

\milcestacion{milcGradeDark}{Ruta MILC: así vamos a aprender}
\begin{tabularx}{\textwidth}{>{\centering\arraybackslash}X>{\centering\arraybackslash}X>{\centering\arraybackslash}X>{\centering\arraybackslash}X}
\phasecard{1}{milcVerde}{milcGris}{Escucha\\[-1mm]\small Observo, escucho y pregunto.} &
\phasecard{2}{milcTurquesa}{milcGris}{Entender\\[-1mm]\small Sistematizo y ordeno.} &
\phasecard{3}{milcMagenta}{milcGris}{Hacer\\[-1mm]\small Construyo y aplico.} &
\phasecard{4}{milcVino}{milcGris}{Cierre\\[-1mm]\small Evalúo y reflexiono.}
\end{tabularx}


\puente{Antes de la teoría vas a listar las mediciones que ya te gobiernan. No las que estudiaste: las que te pasan encima esta semana.}

\milcestacion{milcVerde}{Estación 1 de 3 · Escucha}

\begin{softbox}{milcVerde}{milcGris}{Escena para escuchar}
\textbf{Actividad 1 · IDENTIFICA --- Patrones de medida en tu vida} (15 min · individual). (1) Anota cinco mediciones que regulan tu vida ahora mismo. Por ejemplo la nota de un examen o el tiempo de pantalla. También las libras de comida en casa, los seguidores de una cuenta o las calorías de un producto. (2) Escribe para cada una con qué unidad se mide. (3) Escribe quién definió ese patrón, con nombre o con institución. (4) Escribe qué te pasa a ti según el resultado. (5) Marca la que más te afecta y la que menos entiendes.
\end{softbox}

{\sffamily\bfseries\fontsize{8}{10}\selectfont\textcolor{milcNegro!55}{MARCA CUANDO LO HAYAS HECHO}}
\milccheckrow{Tengo cinco mediciones con su unidad y su resultado práctico.}
\milccheckrow{Escribí quién definió cada patrón, con nombre o institución.}
\milccheckrow{Marqué la que más me afecta y la que menos entiendo.}

\begin{infoband}{milcVerde}


\textbf{Lo que va al cuaderno (Actividad 1 · IDENTIFICA).} \textbf{Título:} «Actividad 1 --- Patrones de medida en mi vida». \textbf{Formato:} tabla de 5 filas y 4 columnas (qué se mide / unidad / quién definió el patrón / qué me pasa según el resultado). \textbf{Extensión:} un tercio de página. \textbf{Sabes que terminaste cuando} ninguna casilla de «quién» dice «la sociedad» o «ellos».
\end{infoband}

\puente{Ya tienes tu lista. Ahora vas a entender por qué una balanza funciona, y por qué el patrón compartido es la parte política del aparato.}

\milcestacion{milcTurquesa}{Estación 2 de 3 · Entender}

\begin{softbox}{milcTurquesa}{milcGris}{Lo que vas a entender}
\textbf{Actividad 2 · EXPLICA --- Las tres partes y el patrón compartido} (20 min · parejas). (1) Con tu pareja, dibujen una balanza y rotulen sus tres partes. (2) Escriban con sus palabras qué pasaría si faltara cada una. (3) Tomen dos mediciones de sus listas y decidan cuál tiene patrón compartido y cuál no. (4) Escriban en dos líneas qué se puede discutir cuando el patrón es compartido y qué no se puede discutir cuando no lo es.
\end{softbox}

\milcpaso{1}{milcTurquesa}{\textbf{Las tres partes.} El punto de apoyo, donde el sistema se equilibra. Los dos brazos iguales, uno para lo desconocido y otro para el patrón. Y el patrón conocido, que es lo único que convierte una comparación en una medida.}
\milcpaso{2}{milcTurquesa}{\textbf{El patrón es la parte política.} Si tú y yo aceptamos la misma pesa, podemos discutir el resultado sin pelear. Si cada uno trae la suya, no hay medición posible: hay dos opiniones con aparato.}
\milcpaso{3}{milcTurquesa}{\textbf{Medir no es lo mismo que juzgar.} La balanza dice cuánto pesa. No dice si eso está bien. Cuando un número se usa directamente como juicio, alguien saltó un paso y casi siempre le convino saltarlo.}
\milcpaso{4}{milcTurquesa}{\textbf{Cuatro preguntas para cualquier sistema de medición.} Qué se mide. Con qué patrón. Quién definió ese patrón y cuándo. Qué se decide con el resultado. La cuarta es la que casi nunca se hace.}

\begin{drawbox}{milcTurquesa}{La balanza y el patrón, en una mirada}
\textbf{Punto de apoyo:} donde el sistema se equilibra. \\[1mm] \textbf{Dos brazos iguales:} lo desconocido contra el patrón. \\[1mm] \textbf{Patrón conocido:} lo que ambas partes aceptan. \\[1mm] \textbf{Pregunta:} ¿quién definió el patrón y qué se decide con el resultado?
\end{drawbox}

\begin{softbox}{milcMostaza}{milcGris}{Errores comunes y su corrección}
(1) \textbf{Creer que el número es neutral}: alguien eligió qué medir y qué dejar fuera. (2) \textbf{Confundir medir con juzgar}: el peso es un dato, la valoración es otra cosa. (3) \textbf{Responder «la sociedad» a la pregunta de quién definió el patrón}: eso no es una respuesta, es una forma de no averiguarlo. (4) \textbf{Olvidar la unidad}: sin unidad, la cifra no dice nada. (5) \textbf{Medir lo fácil en vez de lo que importa}: pasa todo el tiempo, y casi siempre le conviene a alguien.
\end{softbox}

\begin{infoband}{milcTurquesa}


\textbf{Lo que va al cuaderno (Actividad 2 · EXPLICA).} \textbf{Título:} «Actividad 2 --- Las tres partes y el patrón compartido». \textbf{Formato:} el dibujo de la balanza con sus tres partes rotuladas y qué pasa si falta cada una, más las dos mediciones comparadas. \textbf{Extensión:} media página. \textbf{Sabes que terminaste cuando} puedes decir en voz alta cuál de tus dos mediciones no tiene patrón compartido.
\end{infoband}

\puente{Tienes el mecanismo. Toca construirlo con tus manos y después escribir sobre un sistema de medición que sí te afecta.}

\milcestacion{milcMagenta}{Estación 3 de 3 · Hacer}

\begin{softbox}{milcMagenta}{milcGris}{Cómo vas a construir tu producto}
\textbf{Actividad 3 · CREA --- Balanza casera y reflexión} (30 min · individual). (1) Construye una balanza con una regla o un palo, un hilo y dos vasos o tapas iguales. (2) Ajústala hasta que se equilibre en vacío: si no lo hace, todo lo demás miente. (3) Compara dos objetos y anota cuál pesa más. (4) Dibuja o fotografía la balanza y rotula sus tres partes. (5) Escribe una página sobre una medición de tu lista: qué mide, con qué patrón, quién lo definió, qué decide y a quién le conviene. (6) Intercambia la reflexión con un compañero y señálale si nombró a alguien concreto o se quedó en «la sociedad».
\end{softbox}

\milcpaso{1}{milcMagenta}{\textbf{Tu entrega tiene tres piezas.} La balanza que se equilibra en vacío. El dibujo o la foto con las piezas rotuladas. La página de reflexión sobre un solo sistema de medición.}
\milcpaso{2}{milcMagenta}{\textbf{El equilibrio en vacío es la calibración.} Una balanza que ya está inclinada sin nada encima va a mentir en todas las comparaciones. Ajustar el hilo hasta que quede recta es el paso que casi todos se saltan.}
\milcpaso{3}{milcMagenta}{\textbf{Elige un solo sistema para la reflexión.} Uno bien mirado enseña más que cinco mencionados. Y tiene que ser uno que te afecte de verdad esta semana, no un ejemplo de libro.}
\milcpaso{4}{milcMagenta}{\textbf{Nombra a alguien.} «La sociedad» y «el sistema» no definen patrones. Los definen un ministerio, una empresa, un colegio, una aplicación. Si no sabes cuál, averiguarlo es parte del trabajo.}

\begin{drawbox}{milcMagenta}{Antes de entregar}
\checkbox La balanza se equilibra en vacío. \\[1mm] \checkbox Permite comparar dos objetos y decir cuál pesa más. \\[1mm] \checkbox Hay dibujo o foto con las tres partes rotuladas. \\[1mm] \checkbox La reflexión trata un solo sistema de medición. \\[1mm] \checkbox Nombra quién definió el patrón, con nombre o institución. \\[1mm] \checkbox Dice a quién le conviene que se mida así.
\end{drawbox}

\begin{softbox}{milcMagenta}{milcGris}{Plantilla de trabajo}
\textbf{Balanza.} \textit{Materiales, equilibrio en vacío, comparación hecha.} \\[1mm] \textbf{Dibujo.} \textit{Tres partes rotuladas.} \\[1mm] \textbf{Sistema elegido.} \textit{Qué mide y con qué unidad.} \\[1mm] \textbf{Patrón.} \textit{Quién lo definió y cuándo.} \\[1mm] \textbf{Consecuencia.} \textit{Qué se decide y a quién le conviene.}
\end{softbox}

\begin{infoband}{milcMagenta}


\textbf{Lo que va al cuaderno (Actividad 3 · CREA).} \textbf{Título:} «Actividad 3 --- Balanza casera y reflexión». \textbf{Formato:} el dibujo de la balanza con sus partes rotuladas, la comparación de los dos objetos y la página de reflexión. \textbf{Extensión:} media página más la página de reflexión. \textbf{Sabes que terminaste cuando} tu reflexión nombra a alguien concreto que definió el patrón.
\end{infoband}

\puente{Lo que entregas hoy: la balanza que se equilibra y una página sobre quién define el patrón. La prueba de calidad: tu reflexión nombra a alguien concreto, no «la sociedad».}

\milcestacion{milcMostaza}{Producto de la sesión}
\textbf{Balanza casera que se equilibra + reflexión de una página sobre quién define el patrón}.
\begin{softbox}{milcMostaza}{milcGris}{Criterios de calidad}
(1) La balanza se equilibra en vacío y permite comparar dos objetos. (2) Hay dibujo o foto con las tres partes rotuladas. (3) La reflexión trata un solo sistema de medición que te afecta. (4) La reflexión nombra con nombre o institución a quien definió el patrón. (5) La reflexión dice qué se decide con el resultado y a quién le conviene. (6) Revisaste la reflexión de un compañero y le señalaste si nombró a alguien concreto.
\end{softbox}

\puente{Antes de cerrar, mira lo que hiciste desde cinco lados. Aceptar una medida sin preguntar quién la fijó es aceptar el reparto que viene con ella.}

\milcestacion{milcVino}{Evaluación liberadora · cinco dimensiones}

\begin{softbox}{milcVino}{milcGris}{Sentido de la evaluación}
Evaluar no es solo revisar una nota. Es reconocer qué aprendí, qué cambió en mí, qué emociones manejé, cómo me relaciono con la ciudad y la comunidad, y qué le dejaré a quien viene detrás.
\end{softbox}

\textbf{Mi reflexión liberadora en cinco dimensiones.} Completa cada frase iniciada:

\small
\begin{tabularx}{\textwidth}{>{\bfseries\RaggedRight\arraybackslash}p{3.4cm}Y>{\RaggedRight\arraybackslash}p{4.6cm}}
\rowcolor{milcVino}\color{white}Dimensión & \color{white}\bfseries Mi respuesta (frame starter) & \color{white}\bfseries Algo que descubrí\\
1. Personal & Lo que más cambió en mí fue \linead{6.5cm} & \linead{4.2cm}\\
2. Emocional & La emoción más fuerte fue \linead{2.5cm} y la regulé \linead{3cm} & \linead{4.2cm}\\
3. Ciudadana & En democracia esto importa porque \linead{6cm} & \linead{4.2cm}\\
4. Local (Cartago) & En mi barrio sería útil para \linead{6.5cm} & \linead{4.2cm}\\
5. Intergeneracional & A mi abuela/abuelo le diría \linead{6.5cm} & \linead{4.2cm}\\
\end{tabularx}
\normalsize

\begin{infoband}{milcVino}
\textbf{Para la próxima sustentación.} Vuelve a esta tabla en seis meses. Verás que las respuestas cambian — no porque lo aprendido cambie, sino porque tú cambias. La evaluación liberadora es una conversación contigo a través del tiempo.
\end{infoband}


\puente{Para terminar, tres citas verificadas sobre medir y juzgar: Dussel sobre la persona que no es una cosa, Epicteto sobre no confundir la observación con el juicio, Floridi sobre la mitad de los datos que es basura.}

\milcestacion{milcGradeDark}{Triángulo de pensamiento · cierre filosófico}

\begin{softbox}{milcVino}{milcGris}{Dussel · lente del nosotros}
\textit{«El rostro del hombre se revela como otro cuando se recorta en nuestro sistema de instrumentos como exterior, como alguien, como una libertad que interpela, que provoca, que aparece como el que resiste a la totalización instrumental. No es algo; es alguien.»}

{\footnotesize\itshape\color{milcNegro!70}--- Enrique Dussel · Filosofía de la liberación (1977), §2.4.2.2}

Una balanza pesa cosas y eso está bien. El problema empieza cuando el mismo gesto se aplica a personas y el número reemplaza a alguien. En Saucío la escala agrícola ordenaba vecinos, y hasta los zapatos servían de indicador.

\textbf{¿Qué medición de mi vida me convierte en un número para alguien que no me conoce?}
\end{softbox}
\linead{\linewidth}\\[1mm]
\linead{\linewidth}

\begin{softbox}{milcMostaza}{milcGris}{Estoicismo · Epicteto · Enquiridión, 45 (c. 125 d.C.) · cuidado interior}
\textit{«Cuando ves alguno en el baño que se lava pronto, no digas que se lava mal, sino que se lava muy pronto… En efecto, ¿de dónde aprendiste que hizo mal para formar tal juicio?»}

{\footnotesize\itshape\color{milcNegro!70}--- Epicteto · Enquiridión, 45 (c. 125 d.C.)}

La observación dice qué pasó. El juicio dice si estuvo bien, y es un paso aparte. Casi todos los usos abusivos de una medición consisten en dar ese paso sin decir que se dio.

\textbf{¿Cuándo pasé de decir «esto midió tanto» a decir «esto está mal» sin darme cuenta?}
\end{softbox}
\linead{\linewidth}\\[1mm]
\linead{\linewidth}

\begin{softbox}{milcTurquesa}{milcGris}{Floridi · lente de la infoesfera}
\textit{«Parafraseando un dicho de la publicidad: la mitad de nuestros datos es basura, solo que no sabemos cuál mitad. Lo que necesitamos es entender mejor qué datos vale la pena conservar. (trad. propia)»}

{\footnotesize\itshape\color{milcNegro!70}--- Luciano Floridi · Big data and their epistemological challenge (2012)}

Medir mucho no es medir bien. Una aplicación puede registrar cien cosas de ti y no capturar la única que explicaría tu semana. Elegir qué vale la pena medir es más difícil que medir.

\textbf{De todo lo que se mide sobre mí, ¿qué falta que sí explicaría cómo me va?}
\end{softbox}
\linead{\linewidth}\\[1mm]
\linead{\linewidth}

\begin{infoband}{milcNegro}
\textbf{Nota para el docente.} Las tres son citas textuales y llevan su referencia debajo. Puedes remitir a la obra en clase.
\end{infoband}

\milcestacion{milcVino}{Mi autoevaluación · marca una casilla por fila}

{\small
\setlength{\extrarowheight}{2.2pt}
\begin{tabularx}{\textwidth}{>{\RaggedRight\arraybackslash\bfseries}p{3.0cm}YYY}
\rowcolor{milcVino}\color{white}\bfseries Criterio & \color{white}\bfseries Lo logré & \color{white}\bfseries En proceso & \color{white}\bfseries Necesito apoyo\\[.6mm]
Comprensión del tema & \milcopcion{Explico las ideas con mis palabras.} & \milcopcion{Reconozco las ideas, falta precisión.} & \milcopcion{Necesito repasar lo básico.}\\[.8mm]
\rowcolor{milcGris}Pensamiento computacional & \milcopcion{Usé los pilares al armar mi producto.} & \milcopcion{Usé algunos pilares.} & \milcopcion{Necesito releer la estación 2.}\\[.8mm]
Aplicación y producto & \milcopcion{Mi producto cumple los criterios.} & \milcopcion{Está armado, con ajustes pendientes.} & \milcopcion{Aún está en construcción.}\\[.8mm]
\rowcolor{milcGris}Reflexión 5 dimensiones & \milcopcion{Respondí cada una con honestidad.} & \milcopcion{Respondí algunas con detalle.} & \milcopcion{Mis respuestas son muy generales.}\\[.8mm]
Triángulo de pensamiento & \milcopcion{Conecté las 3 voces con mi trabajo.} & \milcopcion{Conecté al menos una.} & \milcopcion{No las relaciono con mi proceso.}\\[.8mm]
\end{tabularx}}

\vspace{3mm}
\textbf{Hoy entendí que:}\\[1mm]
\linead{\linewidth}\\[1mm]
\linead{\linewidth}

\textbf{Mi compromiso para la próxima sesión es:}\\[1mm]
\linead{\linewidth}\\[1mm]
\linead{\linewidth}

\begin{softbox}{milcMostaza}{milcGris}{Cita de despedida · Marco Aurelio}
\textit{``El obstáculo en el camino se vuelve el camino. Lo que estorba al camino se vuelve camino.''}\\[1mm]
Cada error de hoy, bien observado, se convierte en pieza del camino que pisarás mañana.
\end{softbox}

\small
\begin{tabularx}{\textwidth}{>{\bfseries}p{3.6cm}Y}
\rowcolor{milcGradeDark!12}Nombre completo: & \linead{10cm}\\
Fecha: & \linead{4cm} \quad Curso/grupo: \linead{4cm}\\
Mi firma: & \linead{9cm}\\
\end{tabularx}
\normalsize

\begin{infoband}{milcGradeDark}
\textbf{Para la próxima sesión.} Trae la balanza armada y la reflexión. Averigua además quién define una de las medidas de tu lista, con nombre de institución. En la sesión 4 vas a armar la línea del tiempo de la técnica, y las unidades de medida van a aparecer varias veces.
\end{infoband}

\end{document}
